\section{Introduction}

\subsection{Project Overview}

\texttt{lz\_codec} is a command-line lossless compressor for raw 8-bit grayscale images,
implemented in C++17 and built with GNU Make.
It was written by Martin Pribylina (xpriby19) as a course project for KKO 2025/2026
at FIT VUT Brno.

The tool compresses and decompresses files using the LZSS sliding-window dictionary algorithm,
optionally combined with a first-order pixel difference model and per-block adaptive scan
direction selection.
All information needed for decompression is stored inside the compressed file; no flags
need to be passed to the decompressor beyond the input and output file paths.

\subsection{Assignment Constraints}

The binary must support the following CLI flags:

\begin{center}
\begin{tabular}{>{\ttfamily}l l}
\toprule
\textbf{Flag} & \textbf{Meaning} \\
\midrule
-c          & Compress mode \\
-d          & Decompress mode \\
-i \textit{file} & Input file path \\
-o \textit{file} & Output file path \\
-w \textit{N}    & Image width in pixels (required for compression) \\
-m          & Enable pixel difference model (compression only) \\
-a          & Enable adaptive 16$\times$16 block scanning (compression only) \\
\bottomrule
\end{tabular}
\end{center}

Key constraints from the assignment:
\begin{itemize}
  \item \textbf{Self-contained format}: all decompression parameters (\texttt{-m}, \texttt{-a}, \texttt{-w}) are stored
        in the compressed file header. The decompressor never needs these flags.
  \item \textbf{No external compression libraries}: the LZSS encoder/decoder is implemented from scratch.
  \item \textbf{C++17}: compiled with GCC on \texttt{merlin.fit.vutbr.cz} (x86-64 Linux).
  \item \textbf{Portable integers}: \texttt{uint8\_t}, \texttt{uint16\_t}, \texttt{uint32\_t} used throughout.
\end{itemize}

\subsection{Supported Compression Modes}

The tool supports four combinations of the two optional pre-processing flags:

\begin{center}
\begin{tabular}{l l l}
\toprule
\textbf{Mode} & \textbf{Scan} & \textbf{Model} \\
\midrule
1 (default) & Static horizontal & Off \\
2           & Static horizontal & Delta model on \\
3           & Adaptive 16$\times$16 & Off \\
4           & Adaptive 16$\times$16 & Delta model on \\
\bottomrule
\end{tabular}
\end{center}

In all modes the core LZSS algorithm performs the actual byte-count reduction.
The scan order and delta model are pre-processing steps that rearrange or transform the
pixel data to increase the repetition density that LZSS exploits.
